% Generated draft: requires mathematical and source review.
\documentclass[11pt]{article}
\usepackage[margin=25mm]{geometry}
\usepackage{fontspec}
\setmainfont{Latin Modern Roman}
\usepackage{amsmath,amssymb,mathtools,unicode-math}
\setmathfont{Latin Modern Math}
\usepackage{enumitem,needspace,xurl,hyperref}
\setlength{\parindent}{0pt}
\setlength{\parskip}{5pt}
\setlength{\emergencystretch}{3em}
\newcommand{\N}{\mathbb{N}}
\newcommand{\Z}{\mathbb{Z}}
\newcommand{\Q}{\mathbb{Q}}
\newcommand{\R}{\mathbb{R}}
\newcommand{\C}{\mathbb{C}}
\newcommand{\F}{\mathbb{F}}
\newcommand{\A}{\mathbb{A}}
\newcommand{\PP}{\mathbb{P}}
\newcommand{\E}{\mathbb{E}}
\newcommand{\eps}{\varepsilon}
\newcommand{\dd}{\,\mathrm{d}}
\begin{document}
\section*{UnsolvedMath ID: 701693}
\textbf{Source problem code:} catalogue entry
\subsection{Definitions and mathematical statement}
A group \(H\) is word-hyperbolic if it has a finite generating set whose Cayley graph, with its graph metric, has \(\delta\)-thin geodesic triangles for some \(\delta\geq 0\). Such a group is non-elementary if it is not virtually cyclic, where virtually cyclic means having a cyclic subgroup of finite index. A non-elementary hyperbolic quotient of a group \(G\) therefore means a surjective homomorphism \(G\twoheadrightarrow H\) onto such a group \(H\).

For a connected orientable surface of finite type, let \(g\) denote its genus, \(n\) its number of punctures, and \(b\) its number of boundary components, if boundary is admitted. A common complexity convention is \(\xi(S)=3g-3+n+b\); for surfaces without boundary this becomes \(3g-3+n\). Under a common mapping class group convention, \(\operatorname{MCG}(S)\) consists of isotopy classes of orientation-preserving self-homeomorphisms, allowing punctures to be permuted and fixing boundary pointwise, with isotopies satisfying the same restrictions. These are explicit working conventions: the supplied record does not specify whether boundary is admitted, its complexity convention, or the precise homeomorphism and isotopy restrictions defining its mapping class group.

\textbf{Question (M. Hagen, 2022; described in the source as a well-known problem).} For every connected orientable hyperbolic surface \(S\) of finite type and complexity at least \(2\), does there exist a non-elementary word-hyperbolic group \(H\) and a surjective homomorphism
\[
\operatorname{MCG}(S)\twoheadrightarrow H?
\]
Equivalently, is there a normal subgroup \(N\trianglelefteq\operatorname{MCG}(S)\) for which \(\operatorname{MCG}(S)/N\) is non-elementary and word-hyperbolic? The intended surface class and mapping class group convention require confirmation against the original source; the working conventions above do not resolve those source omissions.

\textbf{Supplied qualifications.} The record reports an affirmative answer for the closed surface of genus \(2\), citing \url{https://arxiv.org/pdf/2005.00567.pdf}. It also reports that the same paper gives an affirmative answer conditional on residual finiteness of certain hyperbolic groups. A group \(K\) is residually finite if, for every \(k\in K\setminus\{1\}\), there is a homomorphism from \(K\) to a finite group whose value at \(k\) is nonidentity. The record does not identify the particular hyperbolic groups or the exact scope of the conditional result, so a precise conditional theorem cannot be extracted here. It additionally points to the closely related braid-group question \(16.108\), whose statement is not supplied. These are supplied source reports, not independently verified mathematical or literature assessments.
\subsection{Short English statement}
The question asks whether every surface in the stated class has a mapping class group that surjects onto a non-elementary word-hyperbolic group. The supplied record labels the general problem open and reports a positive closed genus-\(2\) case and a conditional result involving residual finiteness. Surface and mapping class group conventions, and the exact conditional hypotheses, need confirmation.
\subsection{Sources}
\begin{itemize}
\item Local UnsolvedMath registry, permanent ID \(701693\), exact source code \texttt{catalogue entry}, title \emph{Kourovka Notebook Problem 20.35}; supplied complete source record. The registry labels the problem \emph{open}; this label is not an independent verification of its current status. The supplied record was updated at \texttt{2026-10-16T14:00:00Z}.
\item \emph{The Kourovka Notebook (Unsolved Problems in Group Theory)}, problem \(20.35\), proposed by M. Hagen in \(2022\). The supplied background identifies the 20th issue (2022) and the 21st issue edition (2026), \texttt{arXiv:1401.0300v48}. The supplied collection description names E. I. Khukhro and V. D. Mazurov as editors. These bibliographic details are reproduced from the record without independent source verification.
\item Paper supplied at \url{https://arxiv.org/pdf/2005.00567.pdf}, cited by the record for the affirmative closed genus-\(2\) case and the conditional residual-finiteness result. No paper title, authors, or precise theorem statement are supplied.
\item Related reference supplied in the statement: Kourovka Notebook question \(16.108\), concerning braid groups. Its full formulation and a separate citation are not supplied.
\end{itemize}
\end{document}
